\section{Properties of continuous functions}
\label{sec:properties-continuous}
The extreme value theorem concerns a continuous function on a closed,
bounded interval. The intermediate value theorem concerns a continuous
function on any interval. These conclusions answer different questions:
the first guarantees attained extreme values, while the second fills in
values between two already attained values.

All continuity statements below are relative to the specified domain.
For example, an inverse defined on a half-closed interval is tested only
at points belonging to that interval. Multiplying an inequality by a
negative number reverses its direction; we will use that fact explicitly.

\exercise{18.1}{A maximum after reflection is a minimum before reflection}
\goal Suppose $f$ is continuous on $[a,b]$ and $-f$ assumes its
maximum at $x_0\in[a,b]$. Show that $f$ assumes its minimum there.

\step{Translate the meaning of maximum into an inequality.}
For every $x\in[a,b]$, the maximum assumption means
\[
 (-f)(x)\le(-f)(x_0),\qquad\text{that is,}\qquad
 -f(x)\le-f(x_0).
\]
Multiplying both sides by $-1$ reverses the order and gives
$f(x)\ge f(x_0)$. Equivalently,
\[
 f(x_0)\le f(x)\qquad\text{for every }x\in[a,b].
\]
Because $x_0$ belongs to the domain, the value $f(x_0)$ is
actually attained, rather than merely being a lower bound.
\answer{The displayed inequality is exactly the definition of a
minimum of $f$ at $x_0$. Continuity is used by the theorem to obtain
the maximum of $-f$; this final order reversal needs no additional
continuity argument.}

\exercise{18.2}{Where the closed-interval proof uses the endpoints}
\goal Locate the failure in the proof of Theorem~18.1 when its
domain $[a,b]$ is replaced by $(a,b)$, and explain the consequences.

\step{1. A bounded input sequence still has a convergent subsequence.}
Every sequence $(x_n)$ in $(a,b)$ is bounded, so
Bolzano--Weierstrass still provides a subsequence $x_{n_k}$
converging to a real number $x_0$. The difficulty is its location.
The inequalities $a<x_{n_k}<b$ imply only $a\le x_0\le b$.
They do not imply the strict inequalities $a<x_0<b$.
Thus the subsequence may tend to an omitted endpoint.

\step{2. Identify the missing step in the boundedness argument.}
In the theorem's contradiction argument, the selected inputs satisfy
$\abs{f(x_n)}>n$. With a closed interval, the subsequential limit
$x_0$ lies in the domain, so continuity gives
$f(x_{n_k})\to f(x_0)$, a finite real number. On an open interval,
$x_0$ may be outside the domain. Then $f(x_0)$ is not defined, and
continuity of $f$ on $(a,b)$ gives no statement at $x_0$.
The contradiction is lost precisely at that step.

For a concrete example, $f(x)=1/x$ is continuous on $(0,1)$.
The sequence $x_n=1/(n+1)$ lies in that interval and tends to the
excluded point $0$, while $f(x_n)=n+1$ is unbounded.
There is no value $f(0)$ to which continuity could force these
outputs to converge.

\step{3. The same missing-domain issue affects attainment of the bounds.}
Even if $f$ is bounded, a sequence $(y_n)$ chosen so that
$f(y_n)\to\sup f((a,b))$ may have its input subsequence tend to
an endpoint. One then cannot identify the limiting output as a
value of $f$ at an allowed input.

For example, $f(x)=x$ is continuous and bounded on $(0,1)$, with
infimum $0$ and supremum $1$. It has neither a minimum nor a
maximum there. In fact, for every $x\in(0,1)$,
\[
 0<\frac{x}{2}<x<\frac{x+1}{2}<1.
\]
The points $x/2$ and $(x+1)/2$ give strictly smaller and strictly
larger function values than $f(x)$. The sequences approaching
the bounds tend to the excluded endpoints $0$ and $1$.
\answer{Bolzano--Weierstrass still works, but its limit need not
belong to the open domain. Both the boundedness contradiction and
the extremum-attainment step require continuity at that limit.
An individual function on an open interval may still be bounded
and attain extrema; those properties are no longer guaranteed.}

\exercise{18.3}{A polynomial on a half-open interval}
\goal Use calculus to determine the extreme values of
\[
 f(x)=x^3-6x^2+9x+1\qquad(0\le x<5).
\]
\idea Find the stationary points, inspect the derivative's sign,
and remember that $5$ is not an allowed input.

\step{1. Differentiate and find the interior stationary points.}
The power rule gives
\[
 f'(x)=3x^2-12x+9=3(x^2-4x+3)=3(x-1)(x-3).
\]
Thus the only zeros of $f'$ are $1$ and $3$, both in the
domain. Its signs on the intervening intervals are
\[
 \begin{array}{c|ccc}
 x&(0,1)&(1,3)&(3,5)\\ \hline
 x-1&-&+&+\\
 x-3&-&-&+\\
 f'(x)&+&-&+
 \end{array}
\]
The mean value theorem therefore makes $f$ increasing on $[0,1]$,
decreasing on $[1,3]$, and increasing on $[3,5)$.

\step{2. Evaluate all the candidates and the omitted endpoint value.}
Substitution gives
\[
 \begin{aligned}
 f(0)&=1,\\
 f(1)&=1-6+9+1=5,\\
 f(3)&=27-54+27+1=1,\\
 f(5)&=125-150+45+1=21.
 \end{aligned}
\]
The final line evaluates the polynomial formula at $5$ to find
the endpoint limit. It does not add $5$ to the stated domain.

\step{3. Establish and locate the global minimum.}
The useful factorization
\[
 f(x)=x(x-3)^2+1
\]
shows directly that $f(x)\ge1$ for every $x\in[0,5)$:
the factor $x$ and the square $(x-3)^2$ are nonnegative.
Equality requires $x(x-3)^2=0$, which holds exactly when
$x=0$ or $x=3$. Both points are in the domain.

\step{4. Explain why there is no global maximum.}
On $[0,3]$, the derivative analysis gives $f(x)\le5$.
On $[3,5)$, strict increase toward the omitted endpoint gives
$f(x)<21$. In addition, $x_n=5-1/n$ lies in $[0,5)$ and
tends to $5$, so continuity of the polynomial formula on $\R$
gives $f(x_n)\to21$. Thus $21$ is the supremum but is not attained.

One can also exclude a maximum without using the word supremum.
For $x\in[3,5)$, the point $y=(x+5)/2$ satisfies $x<y<5$,
and strict increase gives $f(y)>f(x)$. For $x\in[0,3]$,
we have $f(x)\le5$, whereas $f(4)=64-96+36+1=5$ and
strict increase on $[3,5)$ gives $f(9/2)>f(4)=5$.
Thus every allowed input has some allowed input with a larger
function value.
\answer{The minimum is $1$, attained exactly at $x=0$ and $x=3$.
There is no maximum. The supremum is $21$, approached as
$x\to5$ from the left. The value $5$ at $x=1$ is a local
maximum, not a global maximum on $[0,5)$.}

\exercise{18.4}{A missing limit point produces an unbounded continuous function}
\goal Suppose $S\subseteq\R$ contains a sequence $(x_n)$
converging to $x_0\notin S$. Construct a continuous function
on $S$ whose values are unbounded.

\step{1. Choose a function that grows near the excluded point.}
Define
\[
 F(x)=\frac{1}{\abs{x-x_0}}\qquad(x\in S).
\]
Every $x\in S$ differs from $x_0$, so the denominator is positive
and $F(x)$ is a well-defined finite real number. The function is
a quotient of continuous functions with a nonzero denominator,
restricted to $S$, and hence is continuous there.

\step{2. Verify continuity directly, including the domain restriction.}
Fix $y\in S$ and let $d=\abs{y-x_0}>0$.
If $x\in S$ and $\abs{x-y}<d/2$, the triangle inequality yields
\[
 \abs{x-x_0}\ge\abs{y-x_0}-\abs{x-y}>d/2.
\]
The reverse triangle inequality then gives
\[
 \begin{aligned}
 \abs{F(x)-F(y)}
 &=\frac{\abs{\abs{y-x_0}-\abs{x-x_0}}}
         {\abs{x-x_0}\abs{y-x_0}}\\
 &\le\frac{\abs{x-y}}{\abs{x-x_0}d}
 \le\frac{2\abs{x-y}}{d^2}.
 \end{aligned}
\]
Given $\varepsilon>0$, choose
\[
 \delta=\min\left\{\frac d2,\frac{\varepsilon d^2}{2}\right\}>0.
\]
For every $x\in S$ with $\abs{x-y}<\delta$, the estimate is
strictly less than $\varepsilon$. Thus $F$ is continuous at
every $y\in S$, in precisely the domain-relative sense.

\step{3. Use the given sequence to disprove boundedness.}
Let $M>0$. Since $x_n\to x_0$, there is an integer $N$ such
that $n>N$ implies $\abs{x_n-x_0}<1/M$. All these distances
are nonzero because $x_n\in S$ and $x_0\notin S$.
Reciprocating the positive inequality gives
\[
 F(x_n)=\frac1{\abs{x_n-x_0}}>M\qquad(n>N).
\]
Thus the function takes values larger than every proposed positive
bound. In fact, $F(x_n)\to+\infty$.
\answer{The function $F(x)=1/\abs{x-x_0}$ is continuous on $S$
and unbounded there. Its singular point lies outside its domain.}

\exercise{18.5}{Intersecting graphs and a fixed point}
\goal Prove (a) that the stated endpoint inequalities force
two continuous graphs to meet, and (b) that Example~1 follows
as a special case.

\step{(a), first: subtract the functions.}
Assume $f$ and $g$ are continuous on $[a,b]$, with
$f(a)\ge g(a)$ and $f(b)\le g(b)$.
Define $h(x)=f(x)-g(x)$ on this interval. The difference of
continuous functions is continuous, and
\[
 h(a)\ge0,\qquad h(b)\le0.
\]
If $h(a)=0$, take $x_0=a$. If $h(b)=0$, take $x_0=b$.
Either choice already gives $f(x_0)=g(x_0)$.

\step{(a), second: use the strict sign change in the remaining case.}
If neither endpoint value is zero, the weak inequalities become
$h(a)>0$ and $h(b)<0$. Thus $0$ is strictly between those values.
The intermediate value theorem gives $x_0\in(a,b)$ with
$h(x_0)=0$. By the definition of $h$, this equation is
$f(x_0)-g(x_0)=0$, or $f(x_0)=g(x_0)$.
If the interval degenerates to $a=b$, the two given inequalities
at that single point already force equality, so that case also
is covered.

\step{(b) Compare a map of the unit interval with the identity.}
In Example~1, a continuous function $f$ maps $[0,1]$ into
$[0,1]$. Let $g(x)=x$ on $[0,1]$, a continuous polynomial.
The image restriction gives
\[
 f(0)\ge0=g(0),\qquad f(1)\le1=g(1).
\]
Part (a) therefore supplies $x_0\in[0,1]$ with
$f(x_0)=g(x_0)=x_0$. This is the fixed-point conclusion of
the example, and it includes possible fixed points at $0$ or $1$.
\answer{The endpoint order reversal forces an intersection.
Comparing a continuous map $[0,1]\to[0,1]$ with $g(x)=x$
recovers Example~1.}

\exercise{18.6}{A fixed point of cosine inside a specified interval}
\goal Prove that $x=\cos x$ for at least one
$x\in(0,\pi/2)$.

\step{1. Move both sides to one continuous function.}
Set $h(x)=x-\cos x$ on $[0,\pi/2]$.
The identity and cosine are continuous, so their difference is
continuous on this closed interval.

\step{2. Compute opposite nonzero signs at the endpoints.}
The familiar trigonometric values give
\[
 h(0)=0-1=-1<0,\qquad
 h(\pi/2)=\pi/2-0=\pi/2>0.
\]
Consequently the target value $0$ is strictly between the
endpoint values. By the intermediate value theorem there is
$x\in(0,\pi/2)$ with $h(x)=0$.
\answer{At this interior point, $h(x)=0$ means $x-\cos x=0$,
which is the required equation $x=\cos x$.}

\exercise{18.7}{An exponential equation with an interior solution}
\goal Prove that $xe^x=2$ for some $x\in(0,1)$.

\step{1. Form a continuous function whose zero solves the equation.}
Let $h(x)=xe^x-2$ on $[0,1]$. The identity and exponential
are continuous, and the product and subtraction rules make
$h$ continuous as well.

\step{2. Check the signs, including the inequality involving $e$.}
At zero, $e^0=1$, so $h(0)=0\cdot1-2=-2<0$.
At one, $h(1)=e-2$. The identity
$e=\sum_{k=0}^{\infty}1/k!$, stated in Section~16,
shows that
\[
 e\ge1+1+\frac12=\frac52>2,
\]
because the omitted terms are nonnegative. Hence $h(1)>0$.

\step{3. Apply the intermediate value theorem with strict endpoint signs.}
There is $x\in(0,1)$ with $h(x)=0$, since zero lies strictly
between $h(0)$ and $h(1)$.
\answer{The resulting equation is $xe^x-2=0$, so $xe^x=2$
at an allowed interior point.}

\exercise{18.8}{Opposite signs force a real zero}
\goal Suppose $f:\R\to\R$ is continuous and
$f(a)f(b)<0$. Prove that $f$ has a zero strictly between
$a$ and $b$.

\step{1. Extract the endpoint information from the product.}
A negative product of two real numbers means one factor is
strictly positive and the other strictly negative. In particular,
neither factor is zero. Also $a\ne b$, since otherwise the
product would be the square $f(a)^2\ge0$.

Put $\alpha=\min\{a,b\}$ and $\beta=\max\{a,b\}$.
Then $\alpha<\beta$, and $f(\alpha)$ and $f(\beta)$ still
are the two values of opposite signs. Their order may be either
positive then negative, or negative then positive; in both cases
zero lies strictly between them.

\step{2. Restrict to the ordered interval and apply the theorem.}
The restriction of $f$ to $[\alpha,\beta]$ is continuous.
The intermediate value theorem provides $c\in(\alpha,\beta)$
with $f(c)=0$.
\answer{The point $c$ is strictly between the original endpoints,
regardless of which of $a$ and $b$ was larger.}

\exercise{18.9}{Every polynomial of odd degree has a real root}
\goal For a polynomial
\[
 f(x)=a_0+a_1x+\cdots+a_mx^m,
 \qquad m\text{ odd},\quad a_m\ne0,
\]
prove that $f(c)=0$ for some real $c$.
\idea At inputs of sufficiently large magnitude, the leading
odd power dominates all the lower powers. Give an explicit
bound rather than assume a limit theorem for functions from
the later Section~20.

\step{1. Normalize the leading coefficient without changing the zeros.}
Since $a_m\ne0$, define
\[
 g(x)=\frac{f(x)}{a_m}
 =x^m+\sum_{j=0}^{m-1}b_jx^j,
 \qquad b_j=\frac{a_j}{a_m}.
\]
The functions $f$ and $g$ have exactly the same zeros:
$g(x)=0$ if and only if $f(x)=a_mg(x)=0$.
Both are polynomials and hence are continuous on $\R$.

\step{2. Bound the lower-degree part at large positive and negative inputs.}
Set
\[
 C=\sum_{j=0}^{m-1}\abs{b_j}\ge0,
 \qquad R=\max\{1,C\}+1.
\]
Thus $R>1$ and $R>C$. For either $x=R$ or $x=-R$, we
have $\abs{x}=R$. Since $0\le j\le m-1$ and $R\ge1$,
$R^j\le R^{m-1}$. Applying the triangle inequality gives
\[
 \abs{\sum_{j=0}^{m-1}b_jx^j}
 \le\sum_{j=0}^{m-1}\abs{b_j}R^j
 \le C R^{m-1}.
\]

\step{3. Use the odd parity of $m$ to obtain opposite signs.}
At $R$, the preceding estimate gives
\[
 g(R)\ge R^m-CR^{m-1}
 =R^{m-1}(R-C)>0.
\]
At $-R$, the leading term is $(-R)^m=-R^m$ because $m$
is odd, so
\[
 g(-R)\le -R^m+CR^{m-1}
 =-R^{m-1}(R-C)<0.
\]
Both strict signs follow from $R>0$ and $R-C>0$.
The argument includes $m=1$: the lower-degree part then is
just the constant $b_0$.

\step{4. Fill in zero between those two values.}
Continuity of $g$ on $[-R,R]$ and its opposite endpoint
signs give a point $c\in(-R,R)$ with $g(c)=0$ by the
intermediate value theorem. Hence $f(c)=a_mg(c)=0$.
\answer{Every real polynomial of odd degree has a real root.
The proof treats every nonzero leading coefficient, including
a negative one, through the normalization $g=f/a_m$.}

\exercise{18.10}{Equal values exactly one unit apart}
\goal Suppose $f$ is continuous on $[0,2]$ and $f(0)=f(2)$.
Find $x,y\in[0,2]$ with $\abs{y-x}=1$ and $f(x)=f(y)$.

\step{1. Compare the values separated by one unit.}
For $t\in[0,1]$, both $t$ and $t+1$ belong to the domain
$[0,2]$. Define
\[
 h(t)=f(t+1)-f(t).
\]
Translation $t\mapsto t+1$ is continuous, so composition
makes $t\mapsto f(t+1)$ continuous on $[0,1]$.
The other term is a restriction of $f$, so $h$ is continuous.

\step{2. Show that its endpoint values are negatives of one another.}
The given equality $f(2)=f(0)$ gives
\[
 \begin{aligned}
 h(0)&=f(1)-f(0),\\
 h(1)&=f(2)-f(1)=f(0)-f(1)=-h(0).
 \end{aligned}
\]
If $h(0)=0$, choose $t_0=0$. If $h(0)\ne0$, these two
endpoint values have opposite signs. The intermediate value
theorem then supplies $t_0\in(0,1)$ with $h(t_0)=0$.
Thus in every case there is $t_0\in[0,1]$ with $h(t_0)=0$.

\step{3. Translate the zero back into the requested pair.}
Take $x=t_0$ and $y=t_0+1$. Then $x\in[0,1]$ and
$y\in[1,2]$, so both belong to $[0,2]$. Moreover,
\[
 \abs{y-x}=\abs{(t_0+1)-t_0}=1,
 \qquad f(y)-f(x)=h(t_0)=0.
\]
\answer{The constructed pair has precisely the required distance
and equal function values.}

\exercise{18.11}{The decreasing versions of the inverse-function theorems}
\goal Prove (a) that negating a strictly increasing function
makes it strictly decreasing, and (b) state and prove the
strictly decreasing versions of Theorems~18.4 and~18.5.

\step{(a) Negation reverses each strict inequality.}
If $f$ is strictly increasing on $I$, then $x_1<x_2$ in $I$
implies $f(x_1)<f(x_2)$. Multiplication by $-1$ gives
$-f(x_1)>-f(x_2)$. This is exactly the definition of strict
decrease of $-f$ on $I$.
The same reasoning shows that negating a strictly decreasing
function makes it strictly increasing.

\step{(b), the decreasing version of Theorem 18.5.}
\emph{Statement.} If $g$ is strictly decreasing on an interval
$J$ and its image $g(J)=I$ is an interval, then $g$ is
continuous on $J$.

\emph{Proof.} Let $h=-g$. Part (a), with the inequalities
reversed, makes $h$ strictly increasing. Its image is
\[
 h(J)=-I=\{-y:y\in I\}.
\]
Reflection preserves the interval property. Indeed, if
$u<v$ belong to $-I$ and $u<w<v$, then
$-v<-w<-u$. The endpoints $-v,-u$ belong to $I$, so the
interval property gives $-w\in I$, or $w\in-I$.
Thus $-I$ is an interval, with any included endpoints and
any infinite endpoints reflected in the usual way.

The increasing version of Theorem~18.5 applies to $h$:
its domain is the interval $J$, it is strictly increasing,
and its image $-I$ is an interval. Therefore $h$ is continuous.
Negation is continuous, so $g=-h$ is continuous by Theorem~17.3.
Explicitly, the relevant errors are equal:
\[
 \abs{g(x)-g(x_0)}
 =\abs{-h(x)+h(x_0)}
 =\abs{h(x)-h(x_0)}.
\]
Thus the very same $\delta$ that works for $h$ at $x_0$ works
for $g$. This proves the decreasing version of Theorem~18.5.

\step{(b), the decreasing version of Theorem 18.4.}
\emph{Statement.} If $f$ is continuous and strictly decreasing
on an interval $I$, then its image $J=f(I)$ is an interval,
its inverse is a function $f^{-1}:J\to I$, and that inverse
is continuous and strictly decreasing on $J$.

\emph{Proof.} The intermediate value theorem and Corollary~18.3
make $J=f(I)$ an interval. Strict decrease makes $f$ one-to-one:
distinct inputs can be ordered, and that order gives distinct
outputs. Thus every $u\in J$ has exactly one preimage in $I$,
which defines $f^{-1}(u)$.

Take $u<v$ in $J$, and set $x=f^{-1}(u)$ and $y=f^{-1}(v)$.
If $x=y$, then $u=f(x)=f(y)=v$, a contradiction.
If $x<y$, strict decrease gives $u=f(x)>f(y)=v$, also a
contradiction. Hence $x>y$, which means
$f^{-1}(u)>f^{-1}(v)$. The inverse is strictly decreasing.

Its domain $J$ is an interval and its image is exactly $I$,
also an interval. The decreasing version of Theorem~18.5 just
proved therefore makes $f^{-1}$ continuous on $J$.

No openness of either image interval has been assumed.
Continuity at an included endpoint is relative to $J$; the
definition tests only $u\in J$. If the domain is a singleton,
strict monotonicity is vacuous and the singleton inverse is
continuous directly, so the statements still hold.
\answer{Both decreasing theorems have the same conclusions as
their increasing counterparts, with strict decrease replacing
strict increase. Reflection and the interval-image condition
supply the continuity argument.}

\exercise{18.12}{The intermediate value property without continuity}
\goal For
\[
 f(x)=\begin{cases}\sin(1/x),&x\ne0,\\0,&x=0,\end{cases}
\]
show (a) discontinuity at zero and (b) the intermediate value
property on all of $\R$.

\step{(a) Give an input sequence whose outputs do not approach $f(0)$.}
Let
\[
 u_n=\frac1{2\pi n+\pi/2}\qquad(n\in\N).
\]
The positive denominators grow without bound, so $u_n\to0$.
But periodicity of sine gives
\[
 f(u_n)=\sin(2\pi n+\pi/2)=1
 \qquad\text{for every }n,
\]
whereas $f(0)=0$. Thus $f(u_n)$ does not tend to $f(0)$,
contradicting the sequential condition for continuity.
This is the same discontinuity established in Exercise~17.10(b).

\step{(b), preparation: state exactly what must be checked.}
Fix any $a<b$, and let $y$ lie strictly between $f(a)$ and
$f(b)$. Every value of $f$ belongs to $[-1,1]$, so the strict
intermediate value $y$ satisfies $-1<y<1$.
We must find $x\in(a,b)$ with $f(x)=y$.

\step{(b), intervals that stay away from zero.}
If $0\notin[a,b]$, the reciprocal function $1/x$ is
continuous on $[a,b]$, and composition with sine makes $f$
continuous there. The ordinary intermediate value theorem
then supplies the required interior point.

\step{(b), intervals containing zero and some positive points.}
Suppose $0\in[a,b]$ and $b>0$.
For $n\in\N$, put
\[
 u_n=\frac1{2\pi n+\pi/2},\qquad
 v_n=\frac1{2\pi n+3\pi/2}.
\]
Both sequences tend to zero, and $0<v_n<u_n$.
Choose $n$ large enough that $u_n<b$.
Since $a\le0$, the entire interval $[v_n,u_n]$ lies in
$(a,b)$ and excludes zero. On its endpoints,
\[
 f(v_n)=\sin(2\pi n+3\pi/2)=-1,
 \qquad f(u_n)=1.
\]
The function is continuous on $[v_n,u_n]$ by the composition
argument. Since $-1<y<1$, the intermediate value theorem
gives $x\in(v_n,u_n)\subseteq(a,b)$ with $f(x)=y$.

\step{(b), the remaining intervals end at zero from the left.}
If $0\in[a,b]$ but $b\le0$, then $b=0$ and $a<0$.
Choose $n$ large enough that $u_n<-a$.
Then
\[
 a<-u_n<-v_n<0=b.
\]
For nonzero inputs $f$ is odd, because
$f(-t)=\sin(-1/t)=-\sin(1/t)=-f(t)$.
Therefore $f(-u_n)=-1$ and $f(-v_n)=1$.
The function is continuous on the zero-free interval
$[-u_n,-v_n]$. The intermediate value theorem again gives
an interior point with value $y$, and that point lies in
$(a,b)$.

These cases cover every possible position of $[a,b]$.
When the endpoint values are equal there is no strictly
intermediate $y$ to check, so that case imposes no additional
condition.
\answer{The function is discontinuous at zero but has the
intermediate value property on $\R$. Continuity implies
that property, while the property alone does not imply
continuity.}
